\documentclass[12pt]{article}
\usepackage[utf8]{inputenc}
\usepackage[T1]{fontenc}
\usepackage{hyperref}
\title{Yangum Mon Language Resources for K-12 Curriculum Development}
\author{Compiled by Claude Sonnet 5.5 (claude-sonnet-5-5) for Andr\'as Kornai}
\date{}
\begin{document}
\maketitle

\section{Introduction}
Yangum Mon (ISO 639-3: ymo; Glottocode yang1296) is a Torricelli language of Sandaun Province, Papua New Guinea. Wikipedia cites about 960 speakers (2000 census), and an ASJP wordlist records 730. Mon is the main variety of the Yangum dialect cluster, and its close neighbour is Ambrak. The searches found no school materials in or about the language. This report lists only what the searches turned up.

\section{Primary Educational Resources}
No primers, textbooks or school programmes for Yangum Mon appeared in the searches. The language is not known to be taught in schools.

\section{Linguistic and Reference Materials}
\begin{itemize}
\item Wikipedia, ``Yangum language'', an overview of the dialect cluster: \url{https://en.wikipedia.org/wiki/Yangum_language}
\item Wikipedia, ``Torricelli languages'', for the family context: \url{https://en.wikipedia.org/wiki/Torricelli_languages}
\item Glottolog entry for ISO code ymo, with classification and bibliography pointers: \url{https://glottolog.org/resource/languoid/iso/ymo}
\item ASJP wordlist for Yangum Mon, a short comparative word list: \url{https://asjp.clld.org/languages/YANGUM_MON}
\end{itemize}

\section{Religious and Community Resources}
None found in the searches. The Joshua Project result concerned Yangum Dey, a related variety with a different ISO code, so it is not cited here.

\section{Digital and Audio Resources}
No audio, video or corpus resources for Yangum Mon appeared in the searches.

\section{Curriculum Development Context}
Little documentation is publicly visible. Any curriculum would need new fieldwork with the community. Next steps would be to contact the Summer Institute of Linguistics or Papua New Guinea university linguists and to check archives such as Paradisec, none of which was verified here. Yangum Gel, a sister variety, is reported to be nearly extinct, which makes documentation of Mon more urgent.
\end{document}